\section{The child graph in base 10}
\label{sec:consequences}

\subsection{Successor sequences}

Following each entry of the table until it dies, we confirm the result of~\cite{dbts} for $b = 10$:
the successor automaton has no cycle, so every base-10 comma sequence is finite. In addition, every
successor path that enters a decade $k \ge 11$ dies within $44$ decades, and every sequence starting
at $c \cdot 10^j$ with $j \ge 11$ within $45$ decades. In base 3 the same computation recovers
Theorem~9.1 of~\cite{comma}, with a maximum lifetime of $6$ decades.

\subsection{The infinite path is unique}

In the child graph, a path may switch at a branch point in decade $k$ with leading digit $d$ to the
start $(d + 1) \cdot 10^k$. The set of exits reachable from an input is therefore also given by the
table, and the entries from which an infinite path continues are the greatest fixed point of the
condition ``some reachable exit is again such an entry''.

\begin{theorem}
  \label{thm:unique}
  For every decade $k \ge 11$ there is exactly one offset $u_k \in E$ such that an infinite path of
  the base-10 child graph passes through $10^k + u_k$. Consequently the child graph has exactly one
  infinite path starting at a root, namely \oeis{A367620}, which starts at $20$.
\end{theorem}

\begin{proof}
  The first statement is a finite computation on the table. Every infinite path meets every decade
  from some point on, so it passes through $10^k + u_k$ for all large $k$. Since every node has a
  unique parent, two paths through the same node share all earlier nodes; hence all infinite paths
  from roots coincide.
\end{proof}

All $46$ offsets occur as $u_k$, each for between $12$ and $26$ of the $924$ decades of a period.
Theorem~\ref{thm:unique} answers the question at the end of Section~11 of~\cite{comma}, where
several candidate continuations remained, and is the base-10 analogue of Theorem~10.1
of~\cite{comma} for base 3; our computation reproduces that theorem, including the alternation
between the smaller and the larger child.

\subsection{The branch choices are periodic}

Following \oeis{A367620} decade by decade, at each branch point we take the smaller child unless no
infinite path continues from it. This decision depends only on the state $(k \bmod 924, u_k)$, so:

\begin{theorem}
  \label{thm:periodic}
  The branch choices \oeis{A399179} of \oeis{A367620} are periodic from decade $11$ on, with period
  $924$ decades. One period contains $314$ branch points, at $171$ of which the path takes the
  larger child. Before decade $11$ there is a single branch point, $1\,9^8\,18 = 19999999918$ in
  decade $10$, where the path takes the smaller child.
\end{theorem}

In particular, \oeis{A399179} is given by a finite automaton, as suggested in its OEIS entry, and
all its terms can be listed: the first term, followed by the $314$ terms of Appendix~\ref{app:choices}
repeated forever. The first $30$ terms agree with the known ones,
\[
  \texttt{001110011101100001100101111110},
\]
and we reach the next branch point $2\,9^{84}\,27$ at term $\approx 10^{84.79}$, matching the
``about $10^{84.8}$'' of~\cite[\S11]{comma}.

\begin{remark}
  Counting the root $20$ as term $1$, we find the first branch point $19999999918$ at term
  $412\,987\,859$, both via the table and by direct step-by-step counting. Section~11
  of~\cite{comma} and \oeis{A367620} (offset $1$) give $a(412\,987\,860) = 19999999918$.
  \todo{Resolve this off-by-one.}
\end{remark}
